\subsection{选题侧重与成员分工}

近六年的A、B、C题考查的能力各有侧重。A题往往要从物理过程落到方程，还得看住
数值误差；B题的要害在变量、约束和可行方案；C题通常先处理数据口径，再谈预测、
评价和决策。分工按任务侧重展开：队长主攻机理与连续模型，组员1负责
规划和优化，组员2负责数据与统计学习。三人共用一张输入输出表。队长交出状态量和
物理边界，组员1据此写目标与约束，组员2维护参数、特征和误差口径。这样一来，结果
从谁的脚本产生、又被谁接着使用，在代码和正文里都能顺着查下去。

一项内容要记入本阶段成果，须满足三项要求：能用自己的话解释为什么选它，能把公式拆成
程序动作，也能说出它在哪些条件下会失效。思维导图中的方法远不止这些，这次先完成
十三个单元的原理与代码入口梳理；尚未动手的分支，后面用真题逐项补齐。

